\documentclass[10pt,a4paper]{amsart}
\usepackage[margin=19mm]{geometry}
\usepackage{amsmath,amssymb,mathtools}
\usepackage[scr=rsfs,cal=euler]{mathalfa}
\usepackage{xfrac}
\usepackage[unicode,hidelinks,bookmarks=false]{hyperref}
\newcommand{\ie}{i.e.\ }
\newcommand{\cf}{cf.\ }
\newcommand{\resp}{respectively\ }
\newcommand{\Subsection}{\subsection}
\providecommand{\nobreakdash}{\nobreak\hbox{-}}
\setlength{\emergencystretch}{2em}
\multlinegap0pt
\usepackage{amssymb}
\newtheorem{proposition}{Proposition}
\title{A Simple Trigonometric Classification of Quintic Roots}
\author{Sawon Pratiher}
\date{Source: arXiv:2603.28352v1 (CC BY 4.0)}
\begin{document}
\maketitle

\noindent\emph{Source and licence.} A Simple Trigonometric Classification of Quintic Roots. Version arXiv:2603.28352v1 (CC BY 4.0), distributed by arXiv under the Creative Commons Attribution 4.0 International licence, http://creativecommons.org/licenses/by/4.0/, as recorded in the arXiv licence metadata for this exact version and shown on the abstract page https://arxiv.org/abs/2603.28352v1. Full licence notice: \texttt{licenses/arxiv-2603.28352-CC-BY-4.0.txt}.\par\smallskip

\noindent\textit{Editorial scope.} Counted unit: the article's proposition prop:fP (source0.tex lines 97--102). The article's complete original proof of it is delivered with it (source0.tex lines 104--106, one \texttt{proof} environment of 3 lines). No mathematics is written, repaired or re-derived: the statement and the proof are the article's own lines, in the article's own notation, and no retained line is substituted or re-flowed.  Retained uncounted as the source-local context the proof needs: lines 75--77; lines 79--81.  Nothing was omitted from any retained span, so the package carries no omission record.  No retained line contains a citation, so the wrapper carries no bibliography and no bibliography entry.  No retained line contains a figure environment, an image-inclusion command or drawing code, so no image is delivered and the package declares no asset dependency.  Editorial formatting only, in addition: an AMS \texttt{amsart} preamble that restates the article's own declarations and macros used by the retained lines, namely the environments \texttt{proposition} on separate counters, together with the standard packages those lines need.  The retained environments print in this document as Proposition 1 in order of appearance, whereas the article prints the same items with the numbering of its own full document.  The complete native source of the article as distributed in the arXiv e-print for this exact version is bundled unchanged as \texttt{sources/197427/source0.tex}.
\begin{equation}\label{eq:matched}
16\cos^5\theta - 20\cos^3\theta + \alpha\cos^2\theta + (\beta + 5)\cos\theta + \gamma = 0,
\end{equation}

\begin{equation}\label{eq:params}
\boxed{\alpha = \frac{16n}{u^3}, \qquad \beta = \frac{16p}{u^4} - 5, \qquad \gamma = \frac{16q}{u^5}.}
\end{equation}

\begin{proposition}\label{prop:fP}
For all $\theta \in \mathbb{R}$,
\begin{equation}\label{eq:fP}
f(\theta) = \frac{16}{u^5}\,P(u\cos\theta).
\end{equation}
\end{proposition}

\begin{proof}
Equation~(\ref{eq:matched}) reads $16\cos^5\!\theta - 20\cos^3\!\theta + \alpha\cos^2\!\theta + (\beta+5)\cos\theta + \gamma = 16\,P(u\cos\theta)/u^5$, which holds by the definition of $\alpha$, $\beta$, $\gamma$ in~(\ref{eq:params}). The Chebyshev identity lets us replace the first two terms plus $5\cos\theta$ with $\cos 5\theta$, giving $f(\theta)$ on the left.
\end{proof}

\end{document}
